\documentclass[11pt]{article}
\usepackage[T1]{fontenc}
\usepackage{amsmath,amssymb,amsthm}
\usepackage[margin=1in]{geometry}
\usepackage[hidelinks]{hyperref}\hypersetup{pdftitle={Exact Exclusion of A114216 Hits from Index 78,493,253 through 100,000,000},pdfauthor={Xiaojian Zeng},pdfsubject={Certified finite computation for OEIS A114216},pdfkeywords={integer sequence, segmented sieve, exact computation, OEIS A114216}}
\usepackage{doi}

\newtheorem{theorem}{Theorem}
\newtheorem{lemma}[theorem]{Lemma}

\title{Exact Exclusion of A114216 Hits from Index\\
  $78{,}493{,}253$ through $100{,}000{,}000$}
\author{Xiaojian Zeng\\
  PBC School of Finance, Tsinghua University, 100190\\
  \texttt{xiaojianzeng@foxmail.com}}
\date{}

\begin{document}
\maketitle

\begin{abstract}
Let $p_n$ be the $n$th prime, let $a_0=0$, and define
$a_n=\operatorname{oddpart}(a_{n-1}+p_n)$.  We prove by exhaustive exact
computation that $a_n\ne1$ for
$78{,}493{,}252<n\le100{,}000{,}000$.  Two separately implemented segmented
prime sieves replay the recurrence from $a_0=0$, record every hit, and agree
on the complete hit list and on 100 periodic checkpoints.  The calculation
resolves this finite interval only; it makes no assertion about any later
index or about the existence of a last occurrence.
\end{abstract}

\section{Statement and context}

For a positive integer $m$, write
\[
 v_2(m)=\max\{e\ge0:2^e\mid m\},\qquad
 \operatorname{oddpart}(m)=\frac{m}{2^{v_2(m)}}.
\]
The OEIS entry A114216 gives the recurrence considered here
\cite{oeisA114216}.  Zeng's unrefereed preprint reported the occurrence at
$78{,}493{,}252$ and reconstructed its prime and preceding recurrence state
with two exact-integer implementations \cite{Zeng2026SSRN7395098}.  The
companion entry A114217 now records this index in its occurrence list
\cite{oeisA114217}, but a displayed list alone is not an exhaustive
certificate for a later interval.  Our contribution is the following finite
statement.

\begin{theorem}\label{thm:main}
Let $p_n$ be the $n$th prime, with $p_1=2$, and define $a_0=0$ and
\[
 a_n=\operatorname{oddpart}(a_{n-1}+p_n)\qquad(n\ge1).
\]
Then
\[
 a_n\ne1
 \qquad\text{for every }78{,}493{,}252<n\le100{,}000{,}000.
\]
\end{theorem}

Since $a_{n-1}+p_n>0$, the recurrence is always defined.  Moreover,
\begin{equation}\label{eq:hit}
 a_n=1\quad\Longleftrightarrow\quad
 a_{n-1}+p_n\text{ is a power of }2.
\end{equation}
Thus each hit can be tested with exact integer arithmetic.

\section{Exhaustive certificate}

Both programs enumerate primes below the exclusive bound
\[
U=2{,}100{,}000{,}000
\]
and stop only after processing the $100{,}000{,}000$th prime.  Segmentation is
a standard memory-saving form of the sieve of Eratosthenes
\cite{baysHudson1977}.  A published large-scale precedent used this same
algorithmic family for exhaustive finite verification and prime-gap
computation \cite{oliveiraSilvaHerzogPardi2014}; that work does not certify the
calculation reported here.  Here the sieve is implemented in two deliberately
different ways.

Path A constructs the base primes by a dense sieve, treats 2 separately, and
sieves odd-only segments stored as NumPy Boolean arrays.  It computes the
odd part of a positive unsigned integer $s$ as $s/(s\mathbin{\&}(-s))$.
Path B constructs the base primes by trial division, sieves consecutive
full-integer segments in scalar compiled loops, and removes factors of 2 by
repeated division.  The implementations use NumPy and Numba
\cite{harris2020numpy,lam2015numba}, but share no executable source code.

\begin{lemma}\label{lem:sieve}
Each path enumerates, in increasing order and without omission, all primes it
processes through $p_{100000000}$.
\end{lemma}

\begin{proof}
For each path the base-prime set is complete through
$\lfloor\sqrt{U-1}\rfloor=45{,}825$: it contains 4,746 primes and ends at
45,823.  In a segment, the progression for a base prime $q$ begins at the
larger of $q^2$ and the first multiple of $q$ in that segment.  Hence no prime
is marked.  Conversely, every composite $m<U$ has a prime divisor
$q\le\sqrt m\le\sqrt{U-1}$ and satisfies $m\ge q^2$, so it is marked in its
unique segment.  The survivors traversed in increasing segment order are
therefore exactly the primes.  Both enumerations reach
\[
 p_{100000000}=2{,}038{,}074{,}743<U.
\]
\end{proof}

Starting from $(n,a_n)=(0,0)$, each enumerated prime is fed once into the
recurrence.  Path A divides by the low bit and Path B repeatedly divides by
2; both operations give exactly $s/2^{v_2(s)}$.  Induction on the enumerated
prime index therefore identifies each computed state with the defined
$a_n$.  Both programs check for unsigned 64-bit overflow before every
addition.  The maximum state encountered was $2{,}037{,}477{,}755$, and all
processed primes are below $U$, so the arithmetic used in the recorded runs
is far within the unsigned 64-bit range.

The hit test is executed after every recurrence step, not merely at
checkpoints.  Independently, both paths produce the complete hit list
\begin{equation}\label{eq:hits}
 1,\ 2,\ 5,\ 12,\ 14,\ 20,\ 75,\ 33{,}900,\ 78{,}493{,}252
\end{equation}
through index $100{,}000{,}000$.  At the final hit they reconstruct
\[
 p_{78493252}=1{,}579{,}691{,}809,\qquad
 a_{78493251}=567{,}791{,}839,
\]
whose sum is $2{,}147{,}483{,}648=2^{31}$.  There is no member of the target
interval in \eqref{eq:hits}, proving Theorem~\ref{thm:main}.

As additional cross-checks, the paths agree at every millionth index, giving
100 checkpoint tuples, and agree on the terminal value
\[
 a_{100000000}=1{,}626{,}692{,}547.
\]
They also apply 64-bit FNV-1a to the byte stream formed by concatenating, at
each step, the little-endian unsigned 64-bit encodings of $(n,p_n,a_n)$
\cite{nollEtAl2026FNV}, and obtain the same final digest.  This digest is used
only as provenance and integrity metadata to detect an accidental mismatch;
it is not treated as a proof.  Reruns with altered and unequal segment sizes
reproduce the same mathematical outputs.

\section{Reproduction and limitations}

From the project directory, the decisive records can be regenerated and
compared by
\begin{verbatim}
python3 evidence/path_a_numpy.py --target 100000000 \
  --output /tmp/a114216_full_a.json
python3 evidence/path_b_numba.py --target 100000000 \
  --output /tmp/a114216_full_b.json
python3 evidence/compare_exact_runs.py \
  /tmp/a114216_full_a.json /tmp/a114216_full_b.json \
  --output /tmp/a114216_full_comparison.json
\end{verbatim}
The recorded build used Python 3.10.12, NumPy 2.2.6, and Numba 0.67.0.  The
comparison requires equality of every complete hit record, every checkpoint,
the maximum and terminal states, and the per-step stream digest; it also
reports the hits in the requested interval separately.

The proof is an exhaustive finite calculation.  It neither establishes nor
suggests that $78{,}493{,}252$ is the last hit, and it provides no information
about indices greater than $100{,}000{,}000$.

\paragraph{Computational and AI assistance.}
AI assistance was used in developing the replay programs, organizing the
correctness argument, and drafting this article.  The conclusion rests on the
disclosed exact algorithms and records, not on an AI confidence assessment.
The source code, complete hit records, checkpoint records, altered-segment
runs, and comparison program are supplied with the project for inspection and
reproduction.

\footnotesize
\bibliographystyle{plainurl}\bibliography{references}

\end{document}
